% GENERATED FILE. Edit canonical lessons, then run npm run generate:books.
% canonical-source-hash: fnv1a64:5f279e94c8a4efc4
% canonical-lessons: RU-C288-krazha, RU-C288-vor, RU-C288-krast, RU-C288-karmannik, RU-C288-vzlom

\chapter{Theft}
\label{ch:ru-theft}
\hlchaptermodality{288}

\begin{chapteropening}
\textbf{By the end of this chapter:} \emph{I can report a theft: a theft, a thief, stealing, a pickpocket and a break-in.}

Complete the last lesson of chapter 288: I can report a theft: a theft, a thief, stealing, a pickpocket and a break-in.
\end{chapteropening}

\section[\texorpdfstring{\ru{кража} (krázha)}{кража (krázha)}]{\ru{кража} (krázha) — a theft}

\label{lesson:RU-C288-krazha}



\begin{quote}
Before the new one: say the Russian for a clue, then the Russian for a witness.
\end{quote}

\subsection*{You'll want to know: \ru{кража}}
\textbf{\ru{кража}} (\emph{krázha}) — \textquotedblleft{}a theft\textquotedblright{}.

\begin{quote}
\ru{В} \ru{магазине} \ru{была} \ru{кража}.
\end{quote}

\begin{quote}
\emph{V magazíne bylá krázha.}
\end{quote}

\begin{quote}
There was a theft at the shop.
\end{quote}

Feminine. \emph{\ru{Кража} \ru{со} \ru{взломом}} — burglary.

\begin{cousinweb}
\emph{\ru{кража}} is inherited from Proto-Slavic \emph{kradja} \textquotedblleft{}a theft\textquotedblright{} (EtymDB), built on \emph{krasti} \textquotedblleft{}to steal\textquotedblright{}.
\end{cousinweb}

\begin{grammarlens}[title={What you've built}]
One more word for telling, explaining and describing at B1.
\end{grammarlens}

\subsection*{Your turn}
\begin{itemize}
\raggedright
  \item \emph{Say it:} \emph{krázha}
  \item \emph{Say it:} \emph{krázha}, once more
  \item \emph{Say it:} a sentence of your own with \emph{krázha}
  \item \emph{Recall:} say the Russian for a clue, then the Russian for a witness, then say \emph{krázha} again
\end{itemize}

\begin{tcolorbox}[breakable,colback=teal!4,colframe=teal!35!black,title={Before you move on}]
What does \ru{кража} mean? (\textquotedblleft{}A theft\textquotedblright{}.) Say it once more.
\end{tcolorbox}
\section[\texorpdfstring{\ru{вор} (vor)}{вор (vor)}]{\ru{вор} (vor) — a thief}

\label{lesson:RU-C288-vor}



\begin{quote}
Before the new one: say the Russian for a witness, then the Russian for a theft.
\end{quote}

\subsection*{You'll want to know: \ru{вор}}
\textbf{\ru{вор}} (\emph{vor}) — \textquotedblleft{}a thief\textquotedblright{}.

\begin{quote}
\ru{Вор} \ru{взял} \ru{сумку} \ru{и} \ru{убежал}.
\end{quote}

\begin{quote}
\emph{Vor vzyal súmku i ubezhál.}
\end{quote}

\begin{quote}
The thief took the bag and ran off.
\end{quote}

Genitive \emph{\ru{вора}}, plural \emph{\ru{воры}}, \emph{\ru{воров}}.

\begin{cousinweb}
\emph{\ru{вор}} is an old Russian word whose earlier sense was \textquotedblleft{}a rogue, a rebel\textquotedblright{}; its origin is not settled, and EtymDB lists no parent.
\end{cousinweb}

\begin{grammarlens}[title={What you've built}]
One more word for telling, explaining and describing at B1.
\end{grammarlens}

\subsection*{Your turn}
\begin{itemize}
\raggedright
  \item \emph{Say it:} \emph{vor}
  \item \emph{Say it:} \emph{vor}, once more
  \item \emph{Say it:} a sentence of your own with \emph{vor}
  \item \emph{Recall:} say the Russian for a witness, then the Russian for a theft, then say \emph{vor} again
\end{itemize}

\begin{tcolorbox}[breakable,colback=teal!4,colframe=teal!35!black,title={Before you move on}]
What does \ru{вор} mean? (\textquotedblleft{}A thief\textquotedblright{}.) Say it once more.
\end{tcolorbox}
\section[\texorpdfstring{\ru{красть} (krast)}{красть (krast)}]{\ru{красть} (krast) — to steal}

\label{lesson:RU-C288-krast}



\begin{quote}
Before the new one: say the Russian for a theft, then the Russian for a thief.
\end{quote}

\subsection*{You'll want to know: \ru{красть}}
\textbf{\ru{красть}} (\emph{krast}) — \textquotedblleft{}to steal\textquotedblright{}.

\begin{quote}
\ru{Кто}-\ru{то} \ru{украл} \ru{мой} \ru{велосипед}.
\end{quote}

\begin{quote}
\emph{Kto-to ukrál moy velosipéd.}
\end{quote}

\begin{quote}
Someone stole my bicycle.
\end{quote}

Imperfective; perfective \emph{\ru{украсть}} (\emph{ukrást}). Present \emph{\ru{краду}}, \emph{\ru{крадёшь}} (\emph{kradyósh}); past \emph{\ru{крал}}, \emph{\ru{крала}}.

\begin{cousinweb}
\emph{\ru{красть}} is inherited from Proto-Slavic \emph{krasti} \textquotedblleft{}to steal\textquotedblright{}, which EtymDB traces to PIE \emph{kreh\textsubscript{2}d(\textsuperscript{h})-}.
\end{cousinweb}

\begin{grammarlens}[title={What you've built}]
One more word for telling, explaining and describing at B1.
\end{grammarlens}

\subsection*{Your turn}
\begin{itemize}
\raggedright
  \item \emph{Say it:} \emph{krast}
  \item \emph{Say it:} \emph{krast}, once more
  \item \emph{Say it:} a sentence of your own with \emph{krast}
  \item \emph{Recall:} say the Russian for a theft, then the Russian for a thief, then say \emph{krast} again
\end{itemize}

\begin{tcolorbox}[breakable,colback=teal!4,colframe=teal!35!black,title={Before you move on}]
What does \ru{красть} mean? (\textquotedblleft{}To steal\textquotedblright{}.) Say it once more.
\end{tcolorbox}
\section[\texorpdfstring{\ru{карманник} (karmánnik)}{карманник (karmánnik)}]{\ru{карманник} (karmánnik) — a pickpocket}

\label{lesson:RU-C288-karmannik}



\begin{quote}
Before the new one: say the Russian for a thief, then the Russian for to steal.
\end{quote}

\subsection*{You'll want to know: \ru{карманник}}
\textbf{\ru{карманник}} (\emph{karmánnik}) — \textquotedblleft{}a pickpocket\textquotedblright{}.

\begin{quote}
\ru{В} \ru{автобусе} \ru{работал} \ru{карманник}.
\end{quote}

\begin{quote}
\emph{V avtóbuse rabótal karmánnik.}
\end{quote}

\begin{quote}
A pickpocket was at work on the bus.
\end{quote}

Genitive \emph{\ru{карманника}}.

\begin{cousinweb}
\emph{\ru{карманник}} is built on \emph{\ru{карман}} \textquotedblleft{}a pocket\textquotedblright{}, which is usually traced to a Turkic word for a purse; EtymDB lists \emph{\ru{карман}} \textquotedblleft{}pocket\textquotedblright{} with no parent.
\end{cousinweb}

\begin{grammarlens}[title={What you've built}]
One more word for telling, explaining and describing at B1.
\end{grammarlens}

\subsection*{Your turn}
\begin{itemize}
\raggedright
  \item \emph{Say it:} \emph{karmánnik}
  \item \emph{Say it:} \emph{karmánnik}, once more
  \item \emph{Say it:} a sentence of your own with \emph{karmánnik}
  \item \emph{Recall:} say the Russian for a thief, then the Russian for to steal, then say \emph{karmánnik} again
\end{itemize}

\begin{tcolorbox}[breakable,colback=teal!4,colframe=teal!35!black,title={Before you move on}]
What does \ru{карманник} mean? (\textquotedblleft{}A pickpocket\textquotedblright{}.) Say it once more.
\end{tcolorbox}
\section[\texorpdfstring{\ru{взлом} (vzlom)}{взлом (vzlom)}]{\ru{взлом} (vzlom) — a break-in, forced entry}

\label{lesson:RU-C288-vzlom}



\begin{quote}
Before the new one: say the Russian for to steal, then the Russian for a pickpocket.
\end{quote}

\subsection*{You'll want to know: \ru{взлом}}
\textbf{\ru{взлом}} (\emph{vzlom}) — \textquotedblleft{}a break-in, forced entry\textquotedblright{}.

\begin{quote}
\ru{Соседи} \ru{не} \ru{заметили} \ru{взлома}.
\end{quote}

\begin{quote}
\emph{Sosédi ne zamétili vzlóma.}
\end{quote}

\begin{quote}
The neighbours did not notice the break-in.
\end{quote}

\emph{\ru{Со} \ru{взломом}} — by force.

\begin{cousinweb}
\emph{\ru{взлом}} is built on \emph{\ru{вз}-} and the root of \emph{\ru{ломать}} \textquotedblleft{}to break\textquotedblright{}, an inherited Slavic verb: a breaking open. EtymDB lists no parent.
\end{cousinweb}

\begin{grammarlens}[title={What you've built}]
One more word for telling, explaining and describing at B1.
\end{grammarlens}

\subsection*{Your turn}
\begin{itemize}
\raggedright
  \item \emph{Say it:} \emph{vzlom}
  \item \emph{Say it:} \emph{vzlom}, once more
  \item \emph{Say it:} a sentence of your own with \emph{vzlom}
  \item \emph{Recall:} say the Russian for to steal, then the Russian for a pickpocket, then say \emph{vzlom} again
\end{itemize}

\begin{tcolorbox}[breakable,colback=teal!4,colframe=teal!35!black,title={Before you move on}]
What does \ru{взлом} mean? (\textquotedblleft{}A break-in, forced entry\textquotedblright{}.) Say it once more.
\end{tcolorbox}
